\chapter{Capítulo 2}

A Veeries é uma empresa inserida no agronegócio. 
Oferece produtos relacionados à análise da produção, fluxo e operações de commodities, bem como fluxo e operações dos suprimentos e infra-estrutura relacionados a este mercado. 

A empresa conta com cerca de 20 funcionários e pode ser subdividida em duas equipes: uma de Analistas de mercado e outra de desenvolvedores, como assim são chamados na organização. 
Próximo ao produto final, os analistas são responsáveis por analisar as informações coletadas, fazer projeções e montar relatórios. 
Estritamente ligados ao processamento da informação, o time de desenvolvedores é composto por engenheiros, cientistas da computação e estudantes de áreas correlatas. 
São responsáveis pela melhoria do processo, desde a obtenção dos dados até a apresentação ao cliente.

ETL, sigla para Extraction, Transform, Load (extração, transformação e carregamento) é o processo de integração de dados provindo de uma ou mais fontes de informação. 
Inclui processos de limpeza, organização, validação e carregamento em um repositório organizado e é util para consolidar e garantir dados brutos.

Os dados da fonte de informação, sobretudo o seu nível de organização, dita o quão complexo vai ser sua extração. 
O grau de dificuldade aumenta conforme diminui a clareza da estrutura da informação. 
Comumente é utilizada a classificação de dados estruturados, dados semi-estruturados e dados não estruturados. 
Dados estruturados (bancos de dados relacionais, planilhas excel por exemplo) costumam trazer pouco esforço de processamento por conter sua estrutura quase que pronta para consumo em ferramentas de análise.  
Dados semi-estruturados não se encaixam em tabelas mas possuem marcadores e metadados (tais como saída de leitura de sensores, JSON, XML, logs de sistemas, e-mails). 
Já os dados não estruturados não possuem formato definido (documentos de texto livre, PDFs, imagens, vídeos) e por isso exigem mais esforço no processamento.

Este projeto de ETL vai lidar com um tipo de dado que se classificaria entre dados não estruturados e semi-estruturados, pois a fonte de informação está disposta em tabelas documentadas em arquivos de texto PDF aonde é possível localizar a posição relativa dos elementos de texto  


Descrição da empresa, processos, layouts, problemas, indicadores existentes, etc.

No Capítulo 2, deve-se fazer uma descrição da empresa, processos, layouts, problemas, etc., ou seja, mais detalhadamente motivar e enquadrar o problema dentro do que proporão (a ser descrito no próximo capítulo).